% Esquema do banco de dados MySQL — incluído em textos.tex por % Esquema do banco de dados MySQL — incluído em textos.tex por % Esquema do banco de dados MySQL — incluído em textos.tex por % Esquema do banco de dados MySQL — incluído em textos.tex por \input{figuras/esquema-banco}
% Requer no preâmbulo de principal.tex: \usetikzlibrary{positioning}
\begin{figure}[htbp]
    \centering
    \caption{Esquema do banco de dados da aplicação}
    \label{fig:esquema-banco}
    \vspace{-0.2cm}
    \renewcommand{\arraystretch}{1.1}
    \begin{tikzpicture}[
        tabela/.style={
            draw=black,
            thick,
            inner sep=0pt,
            font=\scriptsize\ttfamily,
            fill=white
        },
        rel/.style={thick},
        card/.style={font=\scriptsize, inner sep=2pt}
    ]
        \node[tabela, anchor=north] (estab) at (0,0) {%
            \begin{tabular}{l}
                \multicolumn{1}{c}{\textbf{estabelecimento}} \\ \hline
                id\_estabelecimento (PK) \\
                cnpj \\
                cnpj\_basico (FK) \\
                cnae\_subclasse (FK) \\
                id\_carga (FK) \\
                nome\_fantasia \\
                situacao\_cadastral \\
                data\_inicio \\
                logradouro \\
                numero \\
                bairro \\
                cep
            \end{tabular}};

        \node[tabela, anchor=north east] (ativ) at (-2.9,0) {%
            \begin{tabular}{l}
                \multicolumn{1}{c}{\textbf{atividade}} \\ \hline
                cnae\_subclasse (PK) \\
                descricao \\
                classe \\
                grupo \\
                divisao
            \end{tabular}};

        \node[tabela, anchor=north west] (empresa) at (2.9,0) {%
            \begin{tabular}{l}
                \multicolumn{1}{c}{\textbf{empresa}} \\ \hline
                cnpj\_basico (PK) \\
                razao\_social \\
                porte \\
                capital\_social \\
                opcao\_mei
            \end{tabular}};

        \node[tabela, below=0.9cm of estab] (loc) {%
            \begin{tabular}{l}
                \multicolumn{1}{c}{\textbf{localizacao}} \\ \hline
                id\_estabelecimento (PK, FK) \\
                latitude \\
                longitude \\
                nivel\_juncao \\
                nv\_geo\_coord \\
                qtd\_mesma\_coordenada \\
                id\_area (FK)
            \end{tabular}};

        \node[tabela, anchor=north east] (carga) at (ativ.east |- loc.north) {%
            \begin{tabular}{l}
                \multicolumn{1}{c}{\textbf{carga}} \\ \hline
                id\_carga (PK) \\
                data\_fotografia \\
                data\_download \\
                versoes\_bases \\
                regras\_filtro
            \end{tabular}};

        \node[tabela, anchor=north west] (area) at (empresa.west |- loc.north) {%
            \begin{tabular}{l}
                \multicolumn{1}{c}{\textbf{area}} \\ \hline
                id\_area (PK) \\
                nome \\
                tipo \\
                area\_km2 \\
                geometria
            \end{tabular}};

        \draw[rel] (ativ.east) -- node[card, pos=0.15, above] {1} node[card, pos=0.85, above] {N} (ativ.east -| estab.west);
        \draw[rel] (empresa.west) -- node[card, pos=0.15, above] {1} node[card, pos=0.85, above] {N} (empresa.west -| estab.east);
        \draw[rel] (estab.south) -- node[card, pos=0.2, right] {1} node[card, pos=0.8, right] {1} (loc.north);
        \draw[rel] (carga.north) |- node[card, pos=0.15, right] {1} node[card, pos=0.9, above] {N} ([yshift=0.5cm]estab.south west);
        \draw[rel] (area.west) -- node[card, pos=0.15, above] {1} node[card, pos=0.85, above] {N} (area.west -| loc.east);
    \end{tikzpicture}
    \vspace{-0.2cm}
    \\\textbf{\footnotesize Fonte: Elaborado pelo autor.}
\end{figure}

% Requer no preâmbulo de principal.tex: \usetikzlibrary{positioning}
\begin{figure}[htbp]
    \centering
    \caption{Esquema do banco de dados da aplicação}
    \label{fig:esquema-banco}
    \vspace{-0.2cm}
    \renewcommand{\arraystretch}{1.1}
    \begin{tikzpicture}[
        tabela/.style={
            draw=black,
            thick,
            inner sep=0pt,
            font=\scriptsize\ttfamily,
            fill=white
        },
        rel/.style={thick},
        card/.style={font=\scriptsize, inner sep=2pt}
    ]
        \node[tabela, anchor=north] (estab) at (0,0) {%
            \begin{tabular}{l}
                \multicolumn{1}{c}{\textbf{estabelecimento}} \\ \hline
                id\_estabelecimento (PK) \\
                cnpj \\
                cnpj\_basico (FK) \\
                cnae\_subclasse (FK) \\
                id\_carga (FK) \\
                nome\_fantasia \\
                situacao\_cadastral \\
                data\_inicio \\
                logradouro \\
                numero \\
                bairro \\
                cep
            \end{tabular}};

        \node[tabela, anchor=north east] (ativ) at (-2.9,0) {%
            \begin{tabular}{l}
                \multicolumn{1}{c}{\textbf{atividade}} \\ \hline
                cnae\_subclasse (PK) \\
                descricao \\
                classe \\
                grupo \\
                divisao
            \end{tabular}};

        \node[tabela, anchor=north west] (empresa) at (2.9,0) {%
            \begin{tabular}{l}
                \multicolumn{1}{c}{\textbf{empresa}} \\ \hline
                cnpj\_basico (PK) \\
                razao\_social \\
                porte \\
                capital\_social \\
                opcao\_mei
            \end{tabular}};

        \node[tabela, below=0.9cm of estab] (loc) {%
            \begin{tabular}{l}
                \multicolumn{1}{c}{\textbf{localizacao}} \\ \hline
                id\_estabelecimento (PK, FK) \\
                latitude \\
                longitude \\
                nivel\_juncao \\
                nv\_geo\_coord \\
                qtd\_mesma\_coordenada \\
                id\_area (FK)
            \end{tabular}};

        \node[tabela, anchor=north east] (carga) at (ativ.east |- loc.north) {%
            \begin{tabular}{l}
                \multicolumn{1}{c}{\textbf{carga}} \\ \hline
                id\_carga (PK) \\
                data\_fotografia \\
                data\_download \\
                versoes\_bases \\
                regras\_filtro
            \end{tabular}};

        \node[tabela, anchor=north west] (area) at (empresa.west |- loc.north) {%
            \begin{tabular}{l}
                \multicolumn{1}{c}{\textbf{area}} \\ \hline
                id\_area (PK) \\
                nome \\
                tipo \\
                area\_km2 \\
                geometria
            \end{tabular}};

        \draw[rel] (ativ.east) -- node[card, pos=0.15, above] {1} node[card, pos=0.85, above] {N} (ativ.east -| estab.west);
        \draw[rel] (empresa.west) -- node[card, pos=0.15, above] {1} node[card, pos=0.85, above] {N} (empresa.west -| estab.east);
        \draw[rel] (estab.south) -- node[card, pos=0.2, right] {1} node[card, pos=0.8, right] {1} (loc.north);
        \draw[rel] (carga.north) |- node[card, pos=0.15, right] {1} node[card, pos=0.9, above] {N} ([yshift=0.5cm]estab.south west);
        \draw[rel] (area.west) -- node[card, pos=0.15, above] {1} node[card, pos=0.85, above] {N} (area.west -| loc.east);
    \end{tikzpicture}
    \vspace{-0.2cm}
    \\\textbf{\footnotesize Fonte: Elaborado pelo autor.}
\end{figure}

% Requer no preâmbulo de principal.tex: \usetikzlibrary{positioning}
\begin{figure}[htbp]
    \centering
    \caption{Esquema do banco de dados da aplicação}
    \label{fig:esquema-banco}
    \vspace{-0.2cm}
    \renewcommand{\arraystretch}{1.1}
    \begin{tikzpicture}[
        tabela/.style={
            draw=black,
            thick,
            inner sep=0pt,
            font=\scriptsize\ttfamily,
            fill=white
        },
        rel/.style={thick},
        card/.style={font=\scriptsize, inner sep=2pt}
    ]
        \node[tabela, anchor=north] (estab) at (0,0) {%
            \begin{tabular}{l}
                \multicolumn{1}{c}{\textbf{estabelecimento}} \\ \hline
                id\_estabelecimento (PK) \\
                cnpj \\
                cnpj\_basico (FK) \\
                cnae\_subclasse (FK) \\
                id\_carga (FK) \\
                nome\_fantasia \\
                situacao\_cadastral \\
                data\_inicio \\
                logradouro \\
                numero \\
                bairro \\
                cep
            \end{tabular}};

        \node[tabela, anchor=north east] (ativ) at (-2.9,0) {%
            \begin{tabular}{l}
                \multicolumn{1}{c}{\textbf{atividade}} \\ \hline
                cnae\_subclasse (PK) \\
                descricao \\
                classe \\
                grupo \\
                divisao
            \end{tabular}};

        \node[tabela, anchor=north west] (empresa) at (2.9,0) {%
            \begin{tabular}{l}
                \multicolumn{1}{c}{\textbf{empresa}} \\ \hline
                cnpj\_basico (PK) \\
                razao\_social \\
                porte \\
                capital\_social \\
                opcao\_mei
            \end{tabular}};

        \node[tabela, below=0.9cm of estab] (loc) {%
            \begin{tabular}{l}
                \multicolumn{1}{c}{\textbf{localizacao}} \\ \hline
                id\_estabelecimento (PK, FK) \\
                latitude \\
                longitude \\
                nivel\_juncao \\
                nv\_geo\_coord \\
                qtd\_mesma\_coordenada \\
                id\_area (FK)
            \end{tabular}};

        \node[tabela, anchor=north east] (carga) at (ativ.east |- loc.north) {%
            \begin{tabular}{l}
                \multicolumn{1}{c}{\textbf{carga}} \\ \hline
                id\_carga (PK) \\
                data\_fotografia \\
                data\_download \\
                versoes\_bases \\
                regras\_filtro
            \end{tabular}};

        \node[tabela, anchor=north west] (area) at (empresa.west |- loc.north) {%
            \begin{tabular}{l}
                \multicolumn{1}{c}{\textbf{area}} \\ \hline
                id\_area (PK) \\
                nome \\
                tipo \\
                area\_km2 \\
                geometria
            \end{tabular}};

        \draw[rel] (ativ.east) -- node[card, pos=0.15, above] {1} node[card, pos=0.85, above] {N} (ativ.east -| estab.west);
        \draw[rel] (empresa.west) -- node[card, pos=0.15, above] {1} node[card, pos=0.85, above] {N} (empresa.west -| estab.east);
        \draw[rel] (estab.south) -- node[card, pos=0.2, right] {1} node[card, pos=0.8, right] {1} (loc.north);
        \draw[rel] (carga.north) |- node[card, pos=0.15, right] {1} node[card, pos=0.9, above] {N} ([yshift=0.5cm]estab.south west);
        \draw[rel] (area.west) -- node[card, pos=0.15, above] {1} node[card, pos=0.85, above] {N} (area.west -| loc.east);
    \end{tikzpicture}
    \vspace{-0.2cm}
    \\\textbf{\footnotesize Fonte: Elaborado pelo autor.}
\end{figure}

% Requer no preâmbulo de principal.tex: \usetikzlibrary{positioning}
\begin{figure}[htbp]
    \centering
    \caption{Esquema do banco de dados da aplicação}
    \label{fig:esquema-banco}
    \vspace{-0.2cm}
    \renewcommand{\arraystretch}{1.1}
    \begin{tikzpicture}[
        tabela/.style={
            draw=black,
            thick,
            inner sep=0pt,
            font=\scriptsize\ttfamily,
            fill=white
        },
        rel/.style={thick},
        card/.style={font=\scriptsize, inner sep=2pt}
    ]
        \node[tabela, anchor=north] (estab) at (0,0) {%
            \begin{tabular}{l}
                \multicolumn{1}{c}{\textbf{estabelecimento}} \\ \hline
                id\_estabelecimento (PK) \\
                cnpj \\
                cnpj\_basico (FK) \\
                cnae\_subclasse (FK) \\
                id\_carga (FK) \\
                nome\_fantasia \\
                situacao\_cadastral \\
                data\_inicio \\
                logradouro \\
                numero \\
                bairro \\
                cep
            \end{tabular}};

        \node[tabela, anchor=north east] (ativ) at (-2.9,0) {%
            \begin{tabular}{l}
                \multicolumn{1}{c}{\textbf{atividade}} \\ \hline
                cnae\_subclasse (PK) \\
                descricao \\
                classe \\
                grupo \\
                divisao
            \end{tabular}};

        \node[tabela, anchor=north west] (empresa) at (2.9,0) {%
            \begin{tabular}{l}
                \multicolumn{1}{c}{\textbf{empresa}} \\ \hline
                cnpj\_basico (PK) \\
                razao\_social \\
                porte \\
                capital\_social \\
                opcao\_mei
            \end{tabular}};

        \node[tabela, below=0.9cm of estab] (loc) {%
            \begin{tabular}{l}
                \multicolumn{1}{c}{\textbf{localizacao}} \\ \hline
                id\_estabelecimento (PK, FK) \\
                latitude \\
                longitude \\
                nivel\_juncao \\
                nv\_geo\_coord \\
                qtd\_mesma\_coordenada \\
                id\_area (FK)
            \end{tabular}};

        \node[tabela, anchor=north east] (carga) at (ativ.east |- loc.north) {%
            \begin{tabular}{l}
                \multicolumn{1}{c}{\textbf{carga}} \\ \hline
                id\_carga (PK) \\
                data\_fotografia \\
                data\_download \\
                versoes\_bases \\
                regras\_filtro
            \end{tabular}};

        \node[tabela, anchor=north west] (area) at (empresa.west |- loc.north) {%
            \begin{tabular}{l}
                \multicolumn{1}{c}{\textbf{area}} \\ \hline
                id\_area (PK) \\
                nome \\
                tipo \\
                area\_km2 \\
                geometria
            \end{tabular}};

        \draw[rel] (ativ.east) -- node[card, pos=0.15, above] {1} node[card, pos=0.85, above] {N} (ativ.east -| estab.west);
        \draw[rel] (empresa.west) -- node[card, pos=0.15, above] {1} node[card, pos=0.85, above] {N} (empresa.west -| estab.east);
        \draw[rel] (estab.south) -- node[card, pos=0.2, right] {1} node[card, pos=0.8, right] {1} (loc.north);
        \draw[rel] (carga.north) |- node[card, pos=0.15, right] {1} node[card, pos=0.9, above] {N} ([yshift=0.5cm]estab.south west);
        \draw[rel] (area.west) -- node[card, pos=0.15, above] {1} node[card, pos=0.85, above] {N} (area.west -| loc.east);
    \end{tikzpicture}
    \vspace{-0.2cm}
    \\\textbf{\footnotesize Fonte: Elaborado pelo autor.}
\end{figure}
